\section{Phương trình đường thẳng trong không gian}
\subsection{Đề 1}
\Opensolutionfile{ans}[ans/ans-2-B15-TN]
\caulc
\begin{ex}%[2H5N2-2]
Trong không gian $Oxyz$, đường thẳng $d\colon \heva{&x=-1+2t\\ &y=2-t\\ &z=3t}$ có một véc-tơ chỉ phương là
\choice
{$\overrightarrow{u}_1=(2;1;3)$}
{$\overrightarrow{u}_2=(-1;2;0)$}
{$\overrightarrow{u}_3=(-1;2;3)$}
{\True $\overrightarrow{u}_4=(2;-1;3)$}
\loigiai{
Đường thẳng $d$ có một véc-tơ chỉ phương là $\overrightarrow{u}_4=(2;-1;3)$.
}
\end{ex}

\begin{ex}%[2H5N2-2]
Trong không gian $Oxyz$, một véc-tơ chỉ phương của đường thẳng $d\colon \dfrac{x+2}{3}=\dfrac{y+1}{-2}=\dfrac{z-3}{-1}$ là
\choice
{$\overrightarrow{u}_1=(-2;1;-3)$}
{\True $\overrightarrow{u}_2=(-3;2;1)$}
{$\overrightarrow{u}_3=(3;-2;1)$}
{$\overrightarrow{u}_4=(2;1;3)$}
\loigiai{
Đường thẳng $d$ có một véc-tơ chỉ phương là $\overrightarrow{u}=(3;-2;-1)$, suy ra $\overrightarrow{u}_2=(-3;2;1)$ là véc-tơ chỉ phương của $d$.
}
\end{ex}

\begin{ex}%[2H5N2-1]
Trong không gian $Oxyz$, đường thẳng $d\colon \heva{&x=t\\&y=1-t\\&z=2+t}$ đi qua điểm nào sau đây?
\choice
{$K(1;-1;1)$}
{$E(1;1;2)$}
{$H(1;2;0)$}
{\True $F(0;1;2)$}
\loigiai{
Xét điểm $F(0;1;2)$, ta thấy $\heva{&0=t\\&1=1-t\\&2=2+t}\Leftrightarrow t=0$.\\
Như vậy $F\in d$.
}
\end{ex}

\begin{ex}%[2H5N2-1]
Trong không gian $Oxyz$, đường thẳng $d\colon\dfrac{x-2}{3}=\dfrac{y-3}{-2}=\dfrac{z-4}{3}$ đi qua điểm nào dưới đây?
\choice
{$Q(3;-2;3)$}
{$M(-2;-3;-4)$}
{\True $P(2;3;4)$}
{$N(-3;2;-3)$}
\loigiai{
Đường thẳng $d$ đã cho đi qua điểm $P(2;3;4)$.
}
\end{ex}

\begin{ex}%[2H5H2-5]
    Trong không gian $Oxyz$, cho hai mặt phẳng $(P) \colon x+y-z-2=0$ và $(Q) \colon x-y-5=0$. Đường thẳng $(d)$ song song với cả hai mặt phẳng $(P)$ và $(Q)$ có véc-tơ chỉ phương là
\choice
{$\vv{u}= \left(-1;-1;1\right)$}
{$\vv{u}= \left(2;-1;1\right)$}
{$\vv{u}= \left(1;-2;1\right)$}
{\True $\vv{u}= \left(-1;-1;-2\right)$}
\loigiai{
Gọi $\vec{u}$ là một véc-tơ chỉ phương của đường thẳng $(d)$, ta có
\begin{itemize}
\item $\vv{n}_P= \left(1;1;-1\right)$ là véc-tơ pháp tuyến của $(P) \Rightarrow \vv{n}_P \perp \vec{u}$.
\item $\vv{n}_Q= \left(1;-1;0\right)$ là véc-tơ pháp tuyến của $(Q) \Rightarrow \vv{n}_Q \perp \vec{u}$.
\end{itemize}
Do đó, véc-tơ chỉ phương của $(d)$ là $\vv{u}=\left[\vv{n}_P;\vv{n}_Q\right]= \left(-1;-1;-2\right)$.
}
\end{ex}

\begin{ex}%[2H5H2-2]
    Trong không gian $Oxyz$, cho ba điểm không thẳng hàng $A(-1;3;1)$, $B(0;3;2)$ và $C(1;2;2)$. Đường thẳng $(d)$ vuông góc với mặt phẳng $(ABC)$ có véc-tơ chỉ phương là
\choice
{\True $\vv{u}= \left(1;2;-3\right)$}
{$\vv{u}= \left(1;1;1\right)$}
{$\vv{u}= \left(1;-2;3\right)$}
{$\vv{u}= \left(1;1;-1\right)$}
\loigiai{
Ta có $\vv{AB}= \left(1;0;1\right)$ và $\vv{AC}= \left(2;-1;1\right)$.\\
Đường thẳng $(d)$ vuông góc với $(ABC) \Rightarrow (d)$ đồng thời vuông góc với giá của hai véc-tơ $\vv{AB}$ và $\vv{AC}$.\\
Do đó, véc-tơ chỉ phương của $(d)$ là $\vv{u}=\left[\vv{AB};\vv{AC}\right]= \left(1;1;-1\right)$.
}
\end{ex}

\begin{ex}%[2H5H2-3]
Trong không gian $Oxyz$, phương trình đường thẳng đi qua hai điểm $A(1;2;3)$ và $B(5;4;-1)$ là
\choice
{$\dfrac{x-5}{2}=\dfrac{y-4}{1}=\dfrac{z+1}{2}$}
{$\dfrac{x+1}{4}=\dfrac{y+2}{2}=\dfrac{z+3}{-4}$}
{$\dfrac{x-1}{4}=\dfrac{y-2}{2}=\dfrac{z-3}{4}$}
{\True $\dfrac{x-3}{-2}=\dfrac{y-3}{-1}=\dfrac{z-1}{2}$}
\loigiai{
Đường thẳng $AB$ có một véc-tơ chỉ phương là $\overrightarrow{AB}=(4;2;-4)$.\\ 
Do đó, đường thẳng $AB$ nhận véc-tơ $\overrightarrow{u}=(-2;-1;2)$ làm véc-tơ chỉ phương.\\
Trong các đường thẳng đã cho, có đường thẳng $\dfrac{x-3}{-2}=\dfrac{y-3}{-1}=\dfrac{z-1}{2}$ đi qua điểm $A$ và có véc-tơ chỉ phương $\overrightarrow{u}=(-2;-1;2)$.\\
Do đó nó chính là phương trình đường thẳng $AB$.
}
\end{ex}
    
\begin{ex}%[2H5H2-3]
Trong không gian $Oxyz$, phương trình đường thẳng $d$ đi qua $M(4;3;2)$ và song song với trục tung là
\choice
{$\heva{& x=4+t\\ & y=3\\& z=2}$}
{\True $\heva{& x=4\\ & y=3+t\\& z=2}$}
{$\heva{& x=4\\ & y=3\\& z=2+t}$}
{$\heva{& x=4-t\\ & y=3\\& z=2}$}
\loigiai{
Đường thẳng $d$ song song với trục tung nên nhận véc-tơ đơn vị $\vec{j}=(0;1;0)$ làm véc-tơ chỉ phương.\\
Lại có $M\in d$ nên ta có phương trình đường thẳng $d$ là $\heva{& x=4\\ & y=3+t\\& z=2.}$
}
\end{ex}      

\begin{ex}%[2H5H2-3]
Trong không gian $Oxyz$, phương trình đường thẳng $d$ đi qua $A(3;5;7)$ và song song với đường thẳng $d' \colon \dfrac{x-1}{2}=\dfrac{y-2}{3}=\dfrac{z-3}{4}$ là
\choice
{\True $\heva{& x=3+2t\\ & y=5+3t\\& z=7+4t}$}
{$\heva{& x=2+3t\\ & y=3+5t\\& z=4+7t}$}
{$\heva{& x=1+3t\\ & y=2+5t\\& z=3+7t}$}
{$\heva{& x=1+2t\\ & y=2+3t\\& z=3+4t}$}
\loigiai{
Đường thẳng $d$ song song với đường thẳng $d'$ nên nhận véc-tơ chỉ phương $\vec{u}=(2;3;4)$ của  $d'$ làm véc-tơ chỉ phương.\\
Lại có $A\in d$ nên ta có phương trình đường thẳng $d$ là $\heva{& x=3+2t\\ & y=5+3t\\& z=7+4t.}$
}
\end{ex}  

\begin{ex}%[2H5H2-3]
Trong không gian $Oxyz$, cho tam giác $ABC$ với $A(1;2;3)$, $B(-3;5;7)$, $C(-1;-4;-1)$. Phương trình đường thẳng vuông góc với mặt phẳng $(ABC)$ tại trọng tâm $G$ của tam giác $ABC$ là
\choice
{$d\colon \dfrac{x-1}{2}=\dfrac{y+1}{-4}=\dfrac{z+3}{5}$}
{$d\colon \dfrac{x+1}{2}=\dfrac{y-1}{4}=\dfrac{z-3}{5}$}
{$d\colon \dfrac{x-1}{2}=\dfrac{y+1}{4}=\dfrac{z+3}{5}$}
{\True $d\colon \dfrac{x+1}{2}=\dfrac{y-1}{-4}=\dfrac{z-3}{5}$}
\loigiai{
Tọa độ trọng tâm $G$ là
$$\heva{& x_G= \dfrac{x_A+x_B+x_C}{3}\\ & y_G= \dfrac{y_A+y_B+y_C}{3}\\ & z_G= \dfrac{z_A+z_B+z_C}{3}} \Leftrightarrow \heva{& x_G=-1\\ & y_G=1\\ &z_G=3.}$$
Ta có $\overrightarrow{AB}=(-4;3;4)$, $\overrightarrow{AC}=(-2;-6;-4)$. Gọi $d$ là đường thẳng cần tìm.\\
Do $d$ vuông góc với mặt phẳng $(ABC)$ nên có véc-tơ chỉ phương là
$$\left[\overrightarrow{AB};\overrightarrow{AC}\right]=(12;-24;30)=6(2;-4;5).$$
Chọn $\vec{u}=(2;-4;5)$ là véc-tơ chỉ phương và đi qua trọng tâm $G(-1;1;3)$ nên có phương trình
$$\dfrac{x+1}{2}=\dfrac{y-1}{-4}=\dfrac{z-3}{5}.$$
}
\end{ex} 

\begin{ex}%[2H5V2-3]
Trong không gian $Oxyz$, tìm đường thẳng $\Delta$ nằm trong mặt phẳng $(P)\colon x+2y+z-4=0$ và vuông góc với đường thẳng  $(d)\colon \dfrac{x+1}{2}=\dfrac{y}{1}=\dfrac{z+2}{3}$, biết $\Delta$ đi qua điểm $M\left(1;1;1 \right) $.
\choice
{\True $\dfrac{x-1}{5}=\dfrac{y-1}{-1}=\dfrac{z-1}{-3}$}
{$\dfrac{x-1}{5}=\dfrac{y-1}{1}=\dfrac{z-1}{-3}$}
{$\dfrac{x-1}{5}=\dfrac{y+1}{-1}=\dfrac{z-1}{2}$}
{$\dfrac{x+1}{5}=\dfrac{y+3}{-1}=\dfrac{z-1}{3}$}
\loigiai{
Gọi $ \vec{n}_{(P)} $ và $  \vec{u}_{(d)} $ lần lượt là véc-tơ pháp tuyến của mặt phẳng $(P)$  và véc-tơ chỉ phương của đường thẳng $(d)$.\\
Gọi $ \vec{u}_{\Delta}$ là véc-tơ chỉ phương của $\Delta$.\\
Ta có $\heva{&\vec{u}_\Delta\perp \vec{n}_{(P)}\\&\vec{u}_\Delta\perp \vec{u}_{(d)}}\Rightarrow \vec{u}_{\Delta}=\left[\vec{n}_{(P)},\vec{u}_{(d)} \right]=(5;-1;-3)$ là véc-tơ chỉ phương của $\Delta$.\\
Đường thẳng $\Delta$ đi qua điểm $ M\left(1;1;1 \right)  $.\\
Vậy phương trình đường thẳng $\Delta$ là $\dfrac{x-1}{5}=\dfrac{y-1}{-1}=\dfrac{z-1}{-3}$.
}
\end{ex} 

\begin{ex}%[2H5V2-6]
Trong không gian với hệ toạ độ $Oxyz$, cho mặt phẳng $(P) \colon 6x-2y+z-35=0$ và điểm $A(-1;3;6)$. Gọi $A'$ là điểm đối xứng với $A$ qua $(P)$. Độ dài đoạn $OA'$ bằng
\choice
{ $3\sqrt{26}$}
{ $5\sqrt{3}$}
{$\sqrt{46}$}
{\True $\sqrt{186} $}
\loigiai{
\immini{
Gọi $ d $ là đường thẳng qua $ A $ và vuông góc với $ (P) $.\\
Ta có $ \vv{u}_{d}=\vv{n}_{P}=(6;-2;1) $. \\
Khi đó phương trình đường thẳng $d\colon \heva{& x=-1+6t \\ & y=3-2t \\ & z=6+t.} $\\
Gọi $ H=d \cap (P) $. Do $ H \in d \Rightarrow H(-1+6t;3-2t;6+t)\in (P)$.\\
}{
\begin{tikzpicture}[scale=1.2,font=\footnotesize,dot/.style={circle,inner sep=1pt,fill,label={#1},name=#1},
extended line/.style={shorten >=-#1,shorten <=-#1},
extended line/.default=1cm]
\path
(0,0) coordinate (O)
(2.5,0) coordinate (B)
(B)--(O)--([turn]-130:1.1) coordinate (D)
($(D)+(B)-( O)$) coordinate (C)
($(O)!0.5!(C)$) coordinate (H)
($(H)+(0,1)$) coordinate (A)
($(A)!2!(H)$) coordinate (A')
($(H)+(0,1.5)$) coordinate (H1)node[shift={(-30:3mm)}]{$d$}
($(A)!1.4!(H)$) coordinate (H2)
($(A)!2.5!(H)$) coordinate (H3)
;
\draw (H1)--(H) (H3)--(H2);
\draw[dashed](H)--(H3);
\draw[dashed](H)--(A)node[red,pos=0.7,rotate=-45]{$ || $};
\draw[dashed](H)--(A')node[red,pos=0.7,rotate=-45]{$ || $};
\begin{scope}
\clip (D)--(O)--(B);
\draw[fill=cyan!60,opacity=0.7] (O) circle(0.5cm)node[black,shift={(25:3.7mm)}]{$P$};
\end{scope}
\draw (O)--(B)--(C)--(D)--cycle  ;
\foreach \p/\r in {H/0,A/180,A'/180}
\fill (\p) circle (1pt) node[shift={(\r:3mm)}]{$\p$};
\end{tikzpicture}
}
\noindent
Thay tọa độ điểm  $ H $ vào phương trình mặt phẳng $ (P) $, ta được
\begin{equation*}
6(-1+6t) -2 \cdot (3-2t)+(6+t)-35 =0 \Leftrightarrow t=1 \Rightarrow H\left(5;1;7\right).
\end{equation*}
Ta có $ A'(a; b; c) $ là điểm đối xứng của $ A $ qua $ (P) \Rightarrow H$ là trung điểm $ AA' $.\\
Áp dụng công thức tìm tọa độ trung điểm, ta có $ \heva{& 5=\dfrac{-1+a}{2} \\ & 1=\dfrac{3+b}{2} \\ & 7=\dfrac{6+c}{2}} \Leftrightarrow \heva{& a= 11 \\ & b=-1\\ & c=8.}$\\
Vậy tọa độ điểm đối xứng với điểm $ A $ qua mặt phẳng $ (P) $ là $ A'\left(11;-1;8\right) $, do đó $OA'=\sqrt{186}$.
}
\end{ex}

\Closesolutionfile{ans}
\indapan{6}{ans/ans-2-B15-TN}

\cauds
\Opensolutionfile{ans}[ans/ans-2-B15-TF]

\begin{ex}%[2H5H2-5]
Trong không gian với hệ tọa độ $Oxyz$, cho mặt phẳng $(P)\colon 2x+y+3z+1=0$ và đường thẳng $d\colon \heva{&x=-3+t \\& y=2-2 t \\& z=1}$. Xét tính đúng - sai của các mệnh đề sau.
\choiceTF
{\True Đường thẳng $d$ đi qua điểm $M(-3;2;1)$}
{Đường thẳng $d$ có véc-tơ chỉ phương $\overrightarrow{u}= \left(1;-2;1\right)$}
{Đường thẳng $d$ cắt mặt phẳng $(P)$}
{\True Đường thẳng $d$ nằm trong mặt phẳng $(P)$}
\loigiai{
\begin{itemchoice}
\itemch Đường thẳng $d$ đi qua điểm $M(-3;2;1)$.
\itemch Đường thẳng $d$ có véc-tơ chỉ phương $\overrightarrow{u}= \left(1;-2;0\right)$.
\itemch $(P)$ có véc-tơ pháp tuyến $\vec{n}=(2;1;3)$.\\
Ta có $\vec{u}\cdot \vec{n}=0 \Rightarrow \vec{u}\bot \vec{n} \Rightarrow d$ nằm trong $(P)$ hoặc $d$ song song $(P)$.
\itemch Thế tọa độ $M(-3;2;1)$ vào phương trình mặt phẳng $(P)$
$$2\cdot (-3)+2+3\cdot 1+1=0 \,\, \text{(đúng)}$$
Vậy đường thẳng $d$ nằm trong mặt phẳng $(P)$
\end{itemchoice}
}
\end{ex}

\begin{ex}%[2H5V2-6]
Trong không gian $Oxyz$, cho điểm $A(2;-3;1)$ và đường thẳng $d \colon \dfrac{x+1}{2}=\dfrac{y+2}{-1}=\dfrac{z}{2}$. Xét tính đúng - sai của các mệnh đề sau.
\choiceTF
{Đường thẳng $d$ đi qua điểm $A$}
{\True Đường thẳng $d$ có véc-tơ chỉ phương $\overrightarrow{u}= \left(2;-1;2\right)$}
{\True Hình chiếu vuông góc của điểm $A$ trên đường thẳng $d$ là điểm $H(1;-3;2)$}
{Điểm $A$ cách đường thẳng $d$ một khoảng ngắn nhất bằng $\sqrt{11}$}
\loigiai{
\begin{itemchoice}
\itemch Ta có $\dfrac{2+1}{2} \ne \dfrac{-3+2}{-1} \ne \dfrac{1}{2}$ nên $A \notin d$.
\itemch Đường thẳng $d$ có véc-tơ chỉ phương $\overrightarrow{u}= \left(2;-1;2\right)$.
\itemch Gọi $H$ là hình chiếu vuông góc của điểm $A$ trên đường thẳng $d$.\\
Do $ H \in d \Rightarrow H(-1+2t; -2-t; 2t) \Rightarrow \overrightarrow{AH}= \left(2t-3;-t+1;2t-1\right)$.\\
Vì $\overrightarrow{AH} \perp \overrightarrow{u}$ nên $\overrightarrow{AH} \cdot \overrightarrow{u}=0$
$$\Leftrightarrow 2 \left(2t-3\right)- \left(-t+1\right)+2 \left(2t-1\right)=0 \Leftrightarrow 9t=9 \Leftrightarrow t=1.$$
Vậy hình chiếu vuông góc của $A$ trên đường thẳng $d$ là $H(1;-3;2)$.
\itemch Điểm $A$ cách đường thẳng $d$ một khoảng ngắn nhất bằng $AH=\sqrt2$.
\end{itemchoice}
}
\end{ex}

\begin{ex}%[2H5V2-3]
Trong không gian với hệ toạ độ $Oxyz$, cho mặt phẳng $(P) \colon x+y+z-3=0$ và đường thẳng  $d \colon \dfrac{x}{1}=\dfrac{y+1}{2}=\dfrac{z-2}{-1}$. Xét tính đúng - sai của các mệnh đề sau.
\choiceTF
{Đường thẳng $d$ có véc-tơ chỉ phương là $\overrightarrow{u}= \left(1;2;1\right)$}
{Đường thẳng $d$ vuông góc với mặt phẳng $(P)$}
{\True Đường thẳng $d$ cắt mặt phẳng $(P)$ tại điểm $M(1;1;1)$}
{\True Hình chiếu của $d$ trên $(P)$ có phương trình là $\dfrac{x-1}{1}=\dfrac{y-1}{4}=\dfrac{z-1}{-5}$}
\loigiai{
\begin{itemchoice}
\itemch Đường thẳng $d$ có véc-tơ chỉ phương là $\overrightarrow{u}= \left(1;2;-1\right)$.
\itemch Mặt phẳng $(P)$ có véc-tơ pháp tuyến là $\overrightarrow{n}= \left(1;1;1\right)$.\\
Do đó $\overrightarrow{u}\cdot \overrightarrow{n} \ne 0$ nên đường thẳng $d$ không vuông góc với mặt phẳng $(P)$.
\itemch Xét hệ phương trình
$$\heva{&x=t\\ &y=2t-1\\ &z=-t+2\\ &x+y+z-3=0} \Leftrightarrow \heva{&t=1\\ &x=1\\ &y=1\\ &z=1.}$$
Vậy đường thẳng $d$ cắt mặt phẳng $(P)$ tại điểm $M(1;1;1)$.
\itemch Gọi $(\alpha)$ là mặt phẳng chứa $d$ và vuông góc với mặt phẳng $(P)$.\\
Khi đó $\vv{n}_{\alpha}=[\vv{u}_d;\vv{n}_P]=(3;-2;-1)$. Chọn $M(0;-1;2) \in d \subset (\alpha)$.\\
Phương trình của mặt phẳng $(\alpha)$ là $3x-2(y+1)-(z-2)=0 \Leftrightarrow 3x-2y-z=0$.\\
Hình chiếu vuông góc của $d$ lên $(P)$ là $\Delta=(\alpha) \cap (P)$ nên xét $\heva{&3x-2y-z=0\\&x+y+z-3=0.}$
Chọn $x=1+t$ ta có  $\heva{&3x-2y-z=0\\&x+y+z-3=0} \Leftrightarrow \heva{&x=1+t\\&y=1+4t\\&z=1-5t.}$\\
Vậy phương trình đường thẳng $\Delta$ là $\heva{&x=1+t\\&y=1+4t\\&z=1-5t.}$
\end{itemchoice}
}
\end{ex}

\begin{ex}%[2H5V2-3]
Trong không gian $Oxyz$, cho điểm $M(1 ; 1 ; 4)$, đường thẳng $d_{1}\colon \dfrac{x-10}{7}=\dfrac{y+4}{1}=\dfrac{z-15}{8}$ và đường thẳng $d_{2}\colon \dfrac{x+1}{-3}=\dfrac{y-1}{4}=\dfrac{z}{5}$. Xét tính đúng - sai của các mệnh đề sau.
\choiceTF
{\True Phương trình tham số của đường thẳng $d_2$ là $ \heva{&x=-1-3t\\ &y=1+4t\\&z=5t}$}
{Đường thẳng $d_1$ đi qua điểm $M$}
{Hai đường thẳng $d_1$ và $d_2$ cắt nhau}
{\True Đường thẳng $d$ đi qua $M$, vuông góc với $d_1$ và cắt $d_2$ có phương trình là $\dfrac{x-1}{4}=\dfrac{y-1}{-4}=\dfrac{z-4}{-3}$}
\loigiai{
\begin{itemchoice}
\itemch Ta thấy $d_2$ có véc-tơ chỉ phương là $\overrightarrow{u}_{2}=(-3;4;5)$ và có phương trình tham số là $ \heva{&x=-1-3t\\ &y=1+4t\\&z=5t.}$
\itemch Ta thấy $\dfrac{1-10}{7}\ne\dfrac{1+4}{1}\ne\dfrac{4-15}{8}$ nên $M \notin d_1$.
\itemch Ta có đường thăng $d_1$ đi qua điểm $A(10;-4;15)$ và có véc-tơ chỉ phương $\overrightarrow{u}_1=(7;1;8)$.\\
Lại có đường thăng $d_2$ đi qua điểm $B(-1;1;0)$ và có véc-tơ chỉ phương $\overrightarrow{u}_2=(-3;4;5)$.\\
Do đó $\overrightarrow{AB}= (-11;5;-15)$ và $\left[\overrightarrow{u}_1\cdot\overrightarrow{u}_2\right]\cdot\overrightarrow{AB} \ne 0$ nên hai đường thẳng $d_1$ và $d_2$ không cắt nhau.
\itemch Gọi $N=d\cap d_2$, ta có $N(-1-3t;1+4t;5t)$ và $\overrightarrow{MN}=(-3t-2;4t;5t-4)$. \\
Do $d\perp d_1$\, nên\, $\overrightarrow{MN}\cdot \overrightarrow{u}_{1}=0\Leftrightarrow 7(-3t-2)+4t+8(5t-4)=0\Leftrightarrow t=2$.\\
Suy ra $\overrightarrow{MN}=\left(-8;8;6\right)$ hay $\overrightarrow{u}=(4;-4;-3)$.\\
Vậy $d \colon \dfrac{x-1}{4}=\dfrac{y-1}{-4}=\dfrac{z-4}{-3}$.
\end{itemchoice}
}
\end{ex}

\Closesolutionfile{ans}
\indapan{3}{ans/ans-2-B15-TF}

\Opensolutionfile{ans}[ans/ans-2-B15-TLN]
\caukq

\begin{ex}%[2H5H2-2]
Trong không gian $Oxyz$, cho tam giác $ABC$ với $A(1;2;3)$, $B(-3;5;7)$, $C(-1;-4;-1)$. Đường thẳng vuông góc với mặt phẳng $(ABC)$ có véc-tơ chỉ phương $\overrightarrow{u}= \left(a;-4;b\right)$. Tính $a^2+b^2$.
\shortans{$29$}
\loigiai{
Ta có $\overrightarrow{AB}=(-4;3;4)$, $\overrightarrow{AC}=(-2;-6;-4)$.\\
Gọi $d$ là đường thẳng cần tìm. Do $d$ vuông góc với mặt phẳng $(ABC)$ nên có véc-tơ chỉ phương là
$$[\overrightarrow{AB};\overrightarrow{AC}]=(12;-24;30)=6(2;-4;5).$$
Vậy $\overrightarrow{u}= \left(2;-4;5\right) \Rightarrow \heva{&a=2\\ &b=5} \Rightarrow a^2+b^2=29$.
}
\end{ex} 

\begin{ex}%[2H5H2-5]
Trong không gian với hệ tọa độ $Oxyz$, cho đường thẳng $d \colon \dfrac{x-1}{2}=\dfrac{y-2}{3}=\dfrac{z-3}{4}$ và mặt phẳng $(P)\colon mx+10y+nz-11=0$. Biết rằng mặt phẳng $(P)$ luôn chứa đường thẳng $d$. Tính giá trị của $m+n$.
\shortans{$-21$}
\loigiai{
Đường thẳng $d$ đi qua $M=(1;2;3)$, có véc-tơ chỉ phương $\vec{a}=(2;3;4).$\\
Mặt phẳng $(P)$ có véc-tơ pháp tuyến $\vec{n}=(m;10;n)$.\\
Ta cần $\heva{&\vec{a}\cdot \vec{n}=0\\&M \in (P)}\Leftrightarrow\heva{&2m+30+4n=0\\&m+20+3n-11=0}$
$$\Leftrightarrow\heva{&2m+4n=-30\\&m+3n=-9 } \Leftrightarrow \heva{&m=-27\\ &n=6} \Rightarrow m+n=-21.$$
}
\end{ex}

\begin{ex}%[2H5H2-5]
Trong không gian $Oxyz$, cho đường thẳng $d\colon \dfrac{x+2}{3}=\dfrac{y-5}{-5}=\dfrac{z-2}{-1}$ và mặt phẳng $(P)\colon 2x+z-3=0$. Đường thẳng $d$ cắt mặt phẳng $(P)$ tại $M \left(x_0;y_0;z_0\right)$. Tính $x_0+y_0+z_0$.
\shortans{$2$}
\loigiai{
Phương trình tham số của đường thẳng $d \colon \heva{&x=-2+3t\\ &y=5-5t\\ &z=2-t.}$\\
Xét hệ phương trình
$$\heva{& x=-2+3t\\ & y=5-5t\\ & z=2-t\\ &2x+z-3=0} \Leftrightarrow \heva{&x=-2+3t\\ &y=5-5t\\ &z=2-t\\ &-4+6t+2-t-3=0} \Leftrightarrow \heva{&x=1\\ &y=0\\ &z=1\\ &t=1.}$$
Vậy đường thẳng $d$ cắt mặt phẳng $(P)$ tại $M(1;0;1) \Rightarrow x_0+y_0+z_0=2$.
}
\end{ex}

\begin{ex}%[2H5V2-3]
Trong không gian $Oxyz$, cho đường thẳng $d$ nằm trong $(P)\colon x+2y+z-4=0$, đồng thời $d$ cắt và vuông góc với đường thẳng $d'\colon \heva{&x=-1+2t\\&y=t\\&z=-2+3t}$. Phương trình đường thẳng $d$ có dạng $\dfrac{x-x_0}{a}=\dfrac{y-y_0}{-1}=\dfrac{z-1}{b}$ với $a\cdot b \ne 0$. Tính $\left(x_0+y_0\right)^2+ \left(a-b\right)^2$.
\shortans{$68$}
\loigiai{
Gọi $\vec{u}_d$ là véc-tơ chỉ phương của đường thẳng $d$.\\
Đường thẳng $d'$ có véc-tơ chỉ phương là $\vec{u}_d'= \left(2;1;3\right)$.\\
Mặt phẳng $(P)$ có véc-tơ pháp tuyến là $\vec{n}_P= \left(1;2;1\right)$.\\
Gọi $K=d\cap d'\Rightarrow K\in d'\Rightarrow K(-1+2t;t;-2+3t)$.\\
Vì $d\subset (P)$ nên $K\in (P)$. Do đó  $(-1+2t)+2t+(-2+3t)-4=0\Leftrightarrow t=1\Rightarrow K(1;1;1)$.\\
Ta có $\heva{&\vec{u}_d\perp \vec{n}_{(P)}\\&\vec{u}_d\perp \vec{u}_{d'}} \Rightarrow \vec{u}_d=\left[\vec{n}_{(P)},\vec{u}_{d'}\right]=(5;-1;-3)$ là véc-tơ chỉ phương của $d$.\\
Phương trình $d$ cần tìm là $\dfrac{x-1}{5}=\dfrac{y-1}{-1}=\dfrac{z-1}{-3} \Rightarrow \heva{&x_0=1\\ &y_0=1\\ &a=5\\ &b=-3.}$\\
Vậy $\left(x_0+y_0\right)^2+ \left(a-b\right)^2= \left(1+1\right)^2+ \left(5-(-3)\right)^2=68$.
}
\end{ex} 

\begin{ex}%[2H5V2-3]
Trong không gian $Oxyz$, cho đường thẳng $d$ đi qua điểm $A(2 ;-1 ; 3)$, vuông góc với đường thẳng $d_{1}\colon \dfrac{x}{4}=\dfrac{y-5}{-1}=\dfrac{z+2}{-1}$ và cắt đường thẳng $d_{2}\colon \dfrac{x-1}{2}=\dfrac{y+1}{3}=\dfrac{z-1}{4}$. Phương trình đường thẳng $d$ có dạng $\dfrac{x-x_0}{1}=\dfrac{y+1}{m}=\dfrac{z-z_0}{n}$ với $m\cdot n \ne 0$. Tính $\left(x_0+z_0\right)^2+ \left(m+n\right)^2$.
\shortans{$41$}
\loigiai{
Ta thấy $d_2$ đi qua điểm $M \left(1;-1;1\right)$ và có véc-tơ chỉ phương là $\overrightarrow{u}_{2}=(2;3;4)$ và có phương trình tham số là $\heva{&x=1+2t\\ &y=-1+3t\\&z=1+4t.}$ \\
Gọi $B=d\cap d_2$, ta có $B(1+2t;-1+3t;1+4t)$ và $\overrightarrow{AB}=(2t-1; 3t; 4t-2)$.\\
Do $d\perp d_1$ nên $\overrightarrow{AB}\cdot \overrightarrow{u}_{1}=0\Leftrightarrow 4(2t-1)-3t-(4t-2) =0\Leftrightarrow t-2=0 \Leftrightarrow t=2$.\\
Suy ra $\overrightarrow{AB}=\left(3;6;6\right)$ hay $\overrightarrow{u}=(1;2;2)$.\\
Do đó, phương trình đường thẳng $d \colon \dfrac{x-2}{1}=\dfrac{y+1}{2}=\dfrac{z-3}{2} \Rightarrow \heva{&x_0=2\\ &z_0=3\\ &m=2\\ &n=2.}$\\
Vậy $\left(x_0+z_0\right)^2+ \left(m+n\right)^2= \left(2+3\right)^2+ \left(2+2\right)^2=41$.
}
\end{ex}

\begin{ex}%[2H5V2-3]
    Trong không gian $O x y z$, cho điểm $A(1 ;-1 ; 2)$, mặt phẳng $(P)\colon x+y-2 z+5=0$ và đường thẳng $d\colon \dfrac{x+1}{2}=\dfrac{y}{1}=\dfrac{z-2}{1}$. Đường thẳng $\Delta$ cắt $d$ và $(P)$ lần lượt tại $M$ và $N$ sao cho $A$ là trung điểm của đoạn thẳng $MN$. Phương trình đường thẳng $\Delta$ có dạng $\dfrac{x+9}{a}=\dfrac{y-y_0}{3}=\dfrac{z-z_0}{b}$ với $a\cdot b \ne 0$. Tính $a^2+b^2+y_0^2+z_0^2$.
    \shortans{$136$}
\loigiai{
Điểm $M$ thuộc đường thẳng $d$ nên tọa độ điểm $M$ là $M(-1+2t;t;2+t)$.\\
Điểm $A$ là trung điểm của $MN$ suy ra $N(3-2t;-2-t;2-t)$.\\
Mặt khác điểm $N$ thuộc mặt phẳng $(P)$ nên $3-2t-2-t-2(2-t)+5=0 \Leftrightarrow t=-4$.\\
Từ đó ta có $M(-9;-4;-2)$, $\vv{AM}=(-10;-3;-4)$ là véc-tơ chỉ phương của đường thẳng $\Delta$.\\
Do đó, phương trình đường thẳng $\Delta$ là $\dfrac{x+9}{10}= \dfrac{y+4}{3}= \dfrac{z+2}{4}$.\\
Vậy $a^2+b^2+y_0^2+z_0^2 = 10^2+4^2+(-4)^2+(-2)^2=136$.
}
\end{ex}
\Closesolutionfile{ans}
\indapan{6}{ans/ans-2-B15-TLN}